\part{Part III: 小红书笔记多维度数据分析与探索}
\chapterimage{chapter_head_2.pdf} % Chapter heading image

\chapter{笔记内容质量分析}

\section{文本内容深度与可读性分析}
本研究从文本长度、可读性与情感三个维度考察小红书笔记内容质量与互动表现之间的关系。由于实际数据处理中，可读性指标仅能对部分样本稳定返回结果，本文采用“全量基础清洗 + 抽样计算特征 + 回归稳健性检验”的方式开展分析。

\subsection{变量构建与样本处理}
首先，本文将原始 CSV 数据合并后进行基础清洗。剔除点赞、评论、收藏缺失的样本后，原始样本量为323,598条，清洗后保留297,031条。为降低极端账号对结果的干扰，本文进一步构造互动率指标，并删除粉丝数异常高但互动率过低的疑似买粉样本，最终样本量仍保持为297,031条。

在核心变量构建上，本文将互动表现定义为：
\[
\text{engagement}=\text{likes}+\text{favs}+4\times \text{comments}
\]
并对其取对数得到以缓解长尾分布对回归估计的影响。文本长度方面，本文计算了标题长度、正文长度以及全文长度，其中全文长度作为主分析变量进入回归模型。情感变量则由情感预测值构造得到，进一步计算情绪极端程度
\[
\text{emotion\_extreme} = \text{emotion\_strength} - 0.5
\]
\subsection{文本长度与互动关系}
为了考察文本长度是否影响用户互动，本文首先对全文长度与互动水平进行了可视化分析。结合后续回归结果，长度的一次项系数为负、二次项系数为正，说明文本长度与互动之间存在一定的 U 型关系。在当前样本区间内，较短文本整体上更容易获得较高互动，但当长度继续增加到一定程度后，互动水平可能再次回升。
\begin{figure}[H]
    \centering 
    \includegraphics[width=0.8\textwidth]{clve.pdf}
    \caption{文本长度与互动关系 }
    \label{fig:content_length_engagement}
\end{figure}

\subsection{文本情感与互动关系}
在情感维度上，本文进一步分析了情感强度与互动表现之间的关系。回归结果表明，情感变量与互动水平呈现负相关趋势。情感强度和情绪极端程度在主模型与稳健性模型中均表现为显著负相关，说明在本样本中，越偏离中性、越强烈的情绪表达，平均互动反而更低。
\begin{figure}[H]
    \centering 
    \includegraphics[width=0.8\textwidth]{eve.pdf}
    \caption{文本情感与互动关系 }
    \label{fig:emotion_engagement}
\end{figure}

\subsection{互动分布特征}
从整体分布来看，笔记互动量呈现明显的长尾特征。大多数样本集中在较低互动区间，只有少量内容获得了极高互动。这也解释了为何本文对互动变量采用对数变换后再进行建模，以提高回归估计的稳定性并减弱极端值的影响。
\begin{figure}[H]
    \centering 
    \includegraphics[width=0.8\textwidth]{edl.pdf}
    \caption{互动特征分布 }
    \label{fig:engagement_distribution}
\end{figure}

\subsection{非线性趋势识别}
为了进一步识别长度变量与互动之间的非线性结构，本文对抽样样本进行了 LOWESS 平滑拟合。结果显示，文本长度与互动之间存在较明显的弯曲趋势，这说明简单线性模型不足以完整描述二者关系。因此，在后续 OLS 回归中，本文加入了长度平方项，以刻画可能存在的非线性效应。
\begin{figure}[H]
    \centering 
    \includegraphics[width=0.8\textwidth]{lowess.pdf}
    \caption{非线性趋势识别 }
    \label{fig:length_lowess}
\end{figure}

\subsection{回归结果与稳健性检验}
在主回归模型中，本文以 
$engagement\_log$ 为因变量，以文本长度、长度平方项、情感强度、情绪极端程度和可读性为自变量进行 OLS 估计。结果显示，模型整体显著，但解释力有限，$R^2$约为0.050。这说明文本特征对互动具有统计意义上的解释力，但笔记热度仍受到账号属性、发布时间、封面质量、主题类别等其他因素影响。

具体而言，文本长度一次项系数显著为负，而二次项系数显著为正，说明长度效应具有非线性特征；情感强度与情绪极端程度的系数均显著为负，表明在本样本中，情绪越强或越极端，平均互动反而越低；可读性变量在主模型中边际显著，但在稳健性模型中不再显著，说明其影响较弱，且对样本极端值较为敏感。

进一步地，本文剔除互动前1%的极端值后重新估计模型，发现核心变量的方向和显著性基本保持一致，说明上述结论具有较好的稳健性。换言之，文本长度与情感特征对互动的影响并非由少数极端爆款样本驱动，而是具有较稳定的总体规律。
\subsection{分组分析}
为更直观地观察长度效应，本文进一步将文本长度按分位数分组，比较不同长度区间的平均互动水平。结果表明，文本长度与平均互动之间存在明显的分组差异，整体趋势与回归结果一致，即长度并非简单线性越长越好，而是存在相对更优的区间。
\begin{figure}[H]
    \centering 
    \includegraphics[width=0.8\textwidth]{ol.pdf}
    \caption{分组分析 }
    \label{fig:length_group_analysis}
\end{figure}

\subsection{情感分组结果}
最后，本文对情感强度与情绪极端程度进行了分组比较。结果表明，情感强度较低、情绪更接近中性的样本平均互动更高；而情绪极端程度越高，平均互动水平总体越低。这一结果与回归分析保持一致，进一步说明在小红书笔记场景中，过强或过于极端的情绪表达并不一定带来更高互动。

综合来看，文本内容特征对笔记热度具有显著影响，其中长度呈现非线性关系，情感强度与情绪极端程度总体上与互动负相关，可读性影响相对较弱且不够稳健。该结果说明，小红书用户对内容的响应不仅取决于文本长度本身，也与内容表达方式、情绪呈现方式及其是否符合平台阅读习惯密切相关。

\section{主题相关性与信息价值分析}

在社交媒体中，除了基础的情感与篇幅特征外，内容传递的实际信息量以及“标题-正文”间的一致性，也是决定用户深度互动的关键因素。本节基于 TF-IDF 算法提取文本核心关键词，从“信息密度”与“标题匹配度”两个维度，深入探讨小红书笔记的信息价值分布对互动表现的影响。

\subsection{变量构建与共线性修正}
为量化笔记的信息分布情况，本文首先利用 TF-IDF 算法提取单篇笔记中的高权重核心词汇。“信息密度（info\_density）”定义为正文高权重词汇数占局部总有效词数的比重；“标题匹配度（title\_match）”则定义为标题关键词与正文关键词的交集词数占标题总词数的比例，主要用于衡量笔记是直接“白描”还是存在一定的内容悬念差。

在将上述变量引入回归模型前，考虑到前期引入的文本长度平方项（$content\_len^2$）会与一次项产生较强的自然数学共线性（未标准化模型的条件数曾高达 $2.04 \times 10^6$），本文利用 Z-score 对所有连续型自变量先行进行标准化处理，再由标准化后的长度构造平方项。处理后，各变量的方差膨胀因子（VIF）介于 1.11 至 7.74 之间，模型总体条件数（Cond. No.）降至 8.52。该结果说明数值尺度导致的病态程度明显缓解，但个别变量仍存在中等程度的共线性，解释单个系数时仍需保持谨慎。

\subsection{信息价值回归结果与检验}
本文将标准化后的各特征变量纳入 OLS 回归模型，因变量为对数互动量 $engagement\_log$，而非除以粉丝数后的互动率。回归估计结果如表 \ref{tab:ols_results_info} 所示。整体模型高度显著（F统计量 = 259.7，Prob(F) $< 0.001$）。

\begin{table}[H]
    \centering
    \caption{信息价值与互动表现的多元线性回归结果}
    \label{tab:ols_results_info}
    \begin{tabular}{lcccc}
        \toprule
        \textbf{变量} & \textbf{标准化系数 ($\beta$)} & \textbf{标准误(Std.Err)} & \textbf{t 值} & \textbf{P>|t|} \\
        \midrule
        截距项 (Intercept)     & 5.4472  & 0.021 & 255.55 & $<$0.001 *** \\
        内容长度 ($content\_len$)      & -0.4994 & 0.043 & -11.71 & $<$0.001 *** \\
        内容长度平方 ($content\_len^2$)& 0.0711  & 0.015 & 4.80   & $<$0.001 *** \\
        情感强度 ($emotion\_strength$) & -0.1100 & 0.018 & -5.97  & $<$0.001 *** \\
        极端情感 ($emotion\_extreme$)   & -0.0969 & 0.019 & -5.18  & $<$0.001 *** \\
        信息密度 ($info\_density$)      & 0.1055  & 0.025 & 4.20   & $<$0.001 *** \\
        标题匹配度 ($title\_match$)     & -0.1822 & 0.016 & -11.23 & $<$0.001 *** \\
        \midrule
        \multicolumn{5}{l}{样本量 = 20,000 \quad $R^2$ = 0.072 \quad 总体条件数(Cond. No.) = 8.52} \\
        \bottomrule
    \end{tabular}
    \vspace{1ex}
    \\
    \small{\textit{注：*** 代表 1\% 的显著性水平。}}
\end{table}

由于引入了标准化处理，相较于前述的基础回归模型，本模型中的各项系数代表了不同特征变量在同一量纲下对互动率的相对作用强度。前述关于长度“U型折返”以及极端情绪负面作用的结论在加入新特征后依旧十分稳健，并且本部分揭示了信息维度两个额外的核心机制。

\subsection{信息密度的“干货溢价”效应}
回归结果表明，信息密度与互动表现呈显著正相关（$\beta=0.1055, p<0.001$）。在控制模型中其他文本变量后，高权重词汇占比较高的笔记通常具有更高的平均互动量。这一结果与“信息更集中、更便于检索的内容可能更符合用户决策需求”的解释一致。不过，本文的信息密度来自TF--IDF阈值构造，且模型尚未控制全部账号与主题差异，因此该结果应理解为相关性证据，而非信息密度对互动的确定性因果效应。

\subsection{预期偏差与“悬念留白”机制}
模型同时揭示了一个值得关注的反直觉现象：标题匹配度与互动表现呈显著负相关（$\beta=-0.1822, p<0.001$），且其标准化系数绝对值在模型中仅次于文本长度。这意味着，在当前样本和词汇重合度定义下，标题关键词被正文直接重复得越多，笔记的平均互动量反而越低。

这一现象可能与信息缺口或标题的悬念表达有关：过分直白的重复可能降低读者继续探索的意愿，而适度的信息差可能保留阅读动机。不过，当前指标衡量的是分词后的词汇重合，而不是严格的语义一致性；标题较短、正文缺失或不同主题的语言习惯也可能影响该指标。因此，“悬念留白”只能作为一种可能解释，尚不能由本模型直接验证。

\subsection{信息分布的可视化特征}
为进一步佐证回归模型揭示的两大规律，本文对信息密度、标题匹配度与互动强度的关系进行了二维六边形密度成像拟合（Hexbin Plot）。如图 \ref{fig:info_hexbin} 所示，左侧图中高密度分布带向右侧高互动区的蔓延直观体现了“干货溢价”的底盘支撑效能；而右侧图则清晰地反映了标题匹配度过高时，样本互动数据出现的集中性坍缩，为本节总结的“悬念留白”机制提供了几何维度的实质性支撑。

\begin{figure}[H]
    \centering 
    \includegraphics[width=0.8\textwidth]{hexbin_density.pdf}
    \caption{信息密度与标题匹配度对互动的二维密度影响分析}
    \label{fig:info_hexbin}
\end{figure}

\chapter{笔记形式与视觉呈现分析}

\section{多媒体内容分析}
小红书是一个强视觉化的平台，图片和视频的质量直接影响用户的点击和停留意愿。本研究将对笔记的多媒体内容进行量化分析。

首先，我们将分析图片的基本特征。对于包含图片的笔记，我们将计算其图片数量，并研究图片数量与互动量的关系。通常，一定数量的高质量图片能更好地展示内容，但过多图片可能导致加载缓慢，反而影响体验。我们还将通过图像处理技术提取封面图或首图的关键视觉特征，如：
\begin{itemize}
	\item \textbf{色彩特征：} 主色调、色彩丰富度（饱和度、对比度）。明亮、色彩和谐的图片可能更具吸引力。
	\item \textbf{构图特征：} 是否包含人脸、人脸占比、是否遵循三分法等构图原则。
	\item \textbf{清晰度与美学质量：} 使用图像清晰度算法（如拉普拉斯方差）和预训练的深度学习美学评估模型对图片质量进行评分。
\end{itemize}
我们将通过相关性分析和分组实验（如对比高热度笔记和低热度笔记的图片特征均值）来验证哪些视觉特征与笔记热度显著相关。

其次，对于视频笔记，我们将分析视频时长、封面帧、是否有字幕、背景音乐等特征。视频时长存在一个“甜点区间”，过短可能表达不清，过长则完播率低。我们将通过数据分析寻找不同内容类别下最优的视频时长范围。

\section{排版与格式分析}
在小红书这类强视觉驱动的社交媒体平台上，笔记的排版与格式直接影响用户的阅读体验和信息获取效率。 良好的排版不仅能够降低用户的认知负荷，还能通过视觉引导突出关键信息，从而延长停留时间并激发互动行为。尽管排版特征往往被视为“微内容 ”，但其对传播效果的潜在贡献不可忽视 。本研究旨在系统量化排版与格式特征对笔记热度的影响，为内容创作者提供可操作的优化建议。
\begin{itemize}
	\item \textbf{文本分段：} 正文是否合理分段（通过换行符或分隔符数量衡量），长段落是否使用小标题或序号。
	\item \textbf{重点突出：} 检测正文中是否包含加粗标签（<b>、<strong>） 、特殊符号或表情符号（Unicode Emoji）来强调关键信息。
	\item \textbf{清单使用：} 检测正文是否包含列表结构，或平台特有的列表格式开头的行，是否使用“1. 2. 3.”或“\textbullet{}”等列表形式，使内容条理清晰。特征通过Python的正则表达式库re和emoji库进行提取。
\end{itemize}
这些格式特征虽然看似细微，但可能通过影响阅读流畅度和信息获取效率，间接影响用户的互动意愿。我们将通过构建二元特征（如“是否使用表情符号”）并比较其在不同热度组中的分布，来评估其影响。

\chapterimage{chapter_head_2.pdf}
\chapter{发布者影响力与社交网络分析}

\section{发布者基础影响力分析}
发布者自身的社交资产是笔记获得初始流量的基石。粉丝数量决定了笔记发布时的基础曝光池大小，而历史表现和认证状态则反映了发布者的内容生产能力和用户信任度。本部分旨在量化发布者基础特征对笔记热度的贡献，并识别高粉丝粘性的“优质创作者”特征。   
\begin{figure}[htbp]
    \centering                % 图片居中
    \includegraphics[width=0.8\textwidth]{16.png}  % 宽度设为页面宽度的80%
    \caption{不同粉丝量级笔记的平均点赞数}  % 图片标题
    \label{fig:followers_likes}       % 标签，供引用
\end{figure}
\begin{itemize}
	\item \textbf{粉丝数量（取对数）：} 从用户信息页采集粉丝数，由于分布呈长尾，对其进行自然对数变换（log(粉丝数+1)）以缓解异方差性。计算粉丝数对数与笔记点赞数对数的皮尔逊相关系数，并通过散点图和平滑拟合（LOESS）观察线性关系强度。
	\item \textbf{历史表现：}计算每个发布者在采集窗口前30天内发布的笔记平均点赞数、评论数、收藏数，作为“近期内容质量 ”的代理变量 。若发布者为新用户，则用整体样本均值填充。定义互动率 = 点赞数 / (粉丝数+1)，计算每个笔记的互动率，并分析其分布 。将发布者按互动率分组， 比较不同组别笔记的平均互动量， 以识别“ 高质量 ”发布者（粉丝数不高但互动率高）。
	\item \textbf{认证状态：} 比较认证用户（企业号、专业号、达人）与非认证用户在平均互动量上的差异。认证通常意味着更高的可信度和专业性，可能带来额外的流量权重。 以对数化后的点赞数为因变量，将粉丝数对数、历史平均表现对数、认证状态等作为自变量，构建多元线性回归模型 。 同时，加入粉丝数与内容特征的交互项，探究不同粉丝量级下内容要素的异质性作用。
\end{itemize}

\section{发布者内容垂直度与粉丝画像（模拟）分析}
由于难以直接获取粉丝画像，我们可以通过分析发布者历史笔记的主题分布来推断其内容垂直度。一个高度垂直的发布者（如只发美妆内容）可能拥有更精准、忠诚的粉丝群体，其笔记的互动率和转化率可能更高。我们将计算发布者历史笔记主题分布的熵值，熵值越低代表垂直度越高，并分析垂直度与平均互动率的关系。

\chapterimage{chapter_head_2.pdf}
\chapter{发布时机与上下文分析}

在社交媒体生态中，内容的传播效能不仅取决于文本与视觉质量，还受到发布时境（Temporal Context）与地域文化（Spatial Context）的深刻影响。本章在利用规则引擎剔除了 15,972 条“高粉丝低互动”的疑似买粉样本后，基于 29.6 万条有效数据，从星期节律、中国特色法定节假日以及地域数字经济三个维度，深度剖析小红书笔记传播的时空动力学机制。

\section{时间节律：基于周周期的注意力波动分析}
受限于现代社会工作与学习的周期性规律，用户的线上注意力呈现出明显的周内波动机制。利用单因素方差分析（ANOVA），检验结果表明“星期几效应对数互动量”的主效应极显著（F = 946.56, $p < 0.001$，$\eta^2 = 0.018$）。

数据呈现出两个值得关注的极值节点：第一，\textbf{周一是内容互动的最高峰（平均 log 互动量达 5.95）}。这一反常识现象可用组织行为学中的“心理补偿机制（Psychological Compensation）”来解释——在结束周末回归工作或学业的首日，用户倾向于利用流式媒体进行碎片化的“赛博摸鱼”以缓解职段压力。其二，\textbf{周五是互动的绝对低谷（平均 log 互动量仅为 5.01）}。Tukey HSD 事后检验进一步表明，周五与其他任何工作日的均值差异（Mean Diff 均在 $-0.5$ 以上）均极其显著（$p_{adj} < 0.01$）。这揭示了媒介使用的替代效应（Substitution Effect）：在周末临近前，受众注意力从“线上内容消费”发生了向“线下实体社交”的物理转移。

\begin{figure}[H]
    \centering 
    \includegraphics[width=\textwidth]{weekday_effect.pdf}
    \caption{小红书笔记互动的周周期波动及分指标分布}
    \label{fig:weekday_effect}
\end{figure}

\section{制度性节点冲击：节假日与调休制度的异象}
除常规节律外，国家法定节假日引起的社会作息变动对平台流量池具有强烈的宏观冲击。为了精准刻画这一效应，本文将 2023-2024 年的数据严格划分为：普通工作日、普通周末、节假日以及中国特有的“调休工作日”。

四组样本的 ANOVA 检验显示出极显著的结构差异（F = 305.53, $p < 10^{-190}$）。其中最引人注目的实证发现是：\textbf{调休工作日的互动峰值反常地超越了所有日期类型，平均总互动跃升至 4744（对数互动 6.28），甚至高于真正的法定节假日（4202）。}

一方面，节假日（Welch t = 18.75, $p < 0.001$ 显著高于非节假日）作为大块闲暇时间，确实是深度长决策内容（如旅游攻略、家电选购）消费的集中区。但另一方面，调休工作日的极高数据表征了“强迫性劳动时间”与“注意力涣散”的错位：处于心理疲劳与调休补偿心态下的用户，其移动端浏览频率剧增。在该节点发布的内容，极易借由宏观“摸鱼流量池”脱颖而出。

\begin{figure}[H]
    \centering 
    \includegraphics[width=0.8\textwidth]{holiday_effect.pdf}
    \caption{不同日期类型（含调休特征）对互动表现的异质性冲击}
    \label{fig:holiday_effect}
\end{figure}

\section{空间格局：地域数字文化的非对称扩张}
从数字地理学的视角切入，社交媒体的数据特征往往能够折射出区域经济水平与消费文化的图景。在剔除无效地理标记后，本文梳理了样本分布排名前列的省份（市），并观察到了“规模优势”与“传播效率高地”并非完全重合的现象。

从总量集聚来看，\textbf{广东省以 25,624 条记录确立了断档第一的大盘领跑属性}（超第二名几乎一倍）。这印证了珠三角庞大的电商与图文代运营产业链造就的内容生产渗透率。然而，从互动转化效率来看，其平均对数互动（5.23）却处于腰部水平，体现了过度商业化与内卷带来的流量互相稀释困境。

另一方面，\textbf{江浙沪与四川盆地构筑了“高生产质感+高互动黏性”的传播堡垒}。浙江（5.97）、四川（5.88）与上海（5.87）在拥有巨大基数的同时，维持了全局最高梯队的绝对互动均值。这在宏观上对应了此类核心地带用户对其输出的“精致生活方式（Lifestyle）”以及活跃“松弛社交文化（Chill Culture）”的向外辐射与强认同。

\begin{figure}[H]
    \centering 
    \includegraphics[width=\textwidth]{province_analysis.pdf}
    \caption{小红书内容互动效能的地域热度解构}
    \label{fig:province_analysis}
\end{figure}

\chapterimage{chapter_head_2.pdf}
\chapter{标签策略与话题热度分析}

\section{标签使用策略分析}
话题标签是小红书内容分类、搜索曝光和推荐分发的重要工具。合理使用标签能够将笔记精准推送给目标受众，同时借助热门标签获取泛流量。本部分旨在系统研究标签数量、类型组合及标签热度对笔记传播效果的影响，为创作者提供科学的标签策略指导。
\begin{figure}[htbp]          % h=here, t=top, b=bottom, p=page
    \centering                % 图片居中
    \includegraphics[width=0.8\textwidth]{17.png}  % 宽度设为页面宽度的80%
    \caption{不同标签数量笔记的平均点赞数}  % 图片标题
    \label{fig:tag_count_likes}       % 标签，供引用
\end{figure}
\begin{itemize}
	\item \textbf{标签数量：} 统计笔记使用的话题标签数量，分析其与热度的关系。过多标签可能被视为spam，过少则可能降低曝光机会。统计每篇笔记包含的标签个数，构建连续特征。绘制标签数量的分布直方图，并统计不同标签数量区间下笔记的平均互动量，寻找最优标签数量区间。
	\item \textbf{标签类型组合：}通过预定义的标签词库（基于数据中高频标签的聚类）将标签划分为以下类型：通用热门标签、垂直细分标签、地点标签、品牌标签等等。对每篇笔记，统计各类标签的数量。分析高热度笔记的标签组合策略，例如是否同时使用了“大流量泛标签”和“精准垂类标签”。使用关联规则挖掘或决策树模型，分析高热度笔记中标签类型的共现模式，识别典型的“黄金标签组合 ”。
	\item \textbf{标签热度：} 计算每个标签在数据采集期间出现的总频次，作为其“基础热度”。分析笔记所带标签的平均热度或最高热度是否与笔记自身热度相关。将笔记按标签平均热度分位数分组， 比较各组互动量的差异，并通过回归模型评估在控制粉丝数等因素后标签热度的边际贡献。

\end{itemize}

\section{话题借势与热点响应分析}
社交媒体内容传播具有显著的时效性，紧跟热点事件是获得爆发式流量的常见策略。然而，盲目“蹭热点”可能导致内容与账号定位不符，甚至引发用户反感。本部分旨在量化热点响应对笔记热度的实际增益，并探索最佳介入时机，为内容营销的时效性决策提供数据支撑。
\begin{itemize}
	\item 效果对比：采用倾向得分匹配（PSM）或逆概率加权（IPW）等方法，为每篇“蹭热点”笔记匹配一组特征相似的非热点笔记，比较两组在发布后72小时内的平均互动量差异，从而得到热点响应的净效应。
	\item 时机效应分析：将热点笔记按介入阶段分组，使用方差分析比较各组互动量的均值，并通过回归模型控制其他因素，估计不同阶段的边际收益。同时，引入热点类型（节日型、事件型、娱乐型）与阶段的交互项，探究异质性。
	\item 生命周期建模：对每个热点事件，追踪其相关笔记的互动量随时间衰减曲线，使用生存分析或时序模型拟合热点的影响力衰减速度，为预测热点剩余价值提供依据。
\end{itemize}
对于高度相关的特征对（如相关系数绝对值大于0.8），在后续建模中可能需要考虑只保留其中一个或进行降维处理。

\chapterimage{chapter_head_2.pdf}
\chapter{多维度相关性分析}

\section{特征间相关性分析}
为避免原始文件、筛选文件及中间结果重复计入，本章仅读取文件名中包含“笔记详细信息”的数据，并优先按照笔记链接去重；链接缺失时，采用编号、标题和博主名称的组合去重。完成互动指标清洗后，共得到182,103条可用笔记，其中168,323条具有有效粉丝数。由于现有数据不包含可靠的图片数量、认证状态和发布小时，本章不对这些变量进行推断，而使用实际可获得的粉丝数、文本长度、标签数、标点与表情符号、情感得分、标题正文重合度、星期和地点完整性等特征。

考虑到互动量和粉丝数均呈明显长尾分布，本文分别构造
\[
 y=\log(1+\text{engagement}),\qquad
 x_{fans}=\log(1+\text{fans}),
\]
并同时计算Pearson相关系数与Spearman秩相关系数。前者用于识别线性关系，后者对极端值更稳健，可用于判断变量之间是否存在稳定的单调关系。图\ref{fig:chapter8_pearson}给出了特征间Pearson相关系数矩阵。

\begin{figure}[H]
    \centering
    \includegraphics[width=0.92\textwidth]{chapter8_pearson_heatmap.png}
    \caption{数值特征间的Pearson相关系数矩阵}
    \label{fig:chapter8_pearson}
\end{figure}

结果显示，正文长度与全文长度的相关系数达到$r=0.988$，是唯一满足$|r|\geq0.8$的特征对。这一高度相关主要源于全文长度由标题和正文共同构成，属于变量定义带来的机械相关。在后续多变量模型中不宜同时保留正文长度与全文长度，本文建议保留含义更完整的全文长度。其余变量之间未出现绝对相关系数超过0.8的组合，说明除长度变量外，当前特征集合不存在严重的两两线性共线性。需要强调的是，两两相关较低并不能完全排除多变量共线性，后续回归仍应结合VIF进行诊断。

\section{特征与目标变量相关性分析}
本文首先以对数互动量为连续型目标，分别计算各特征的Pearson和Spearman相关系数，并采用Benjamini--Hochberg方法对多重检验的$p$值进行FDR校正。图\ref{fig:chapter8_target}按照Spearman相关系数绝对值展示变量排序。

\begin{figure}[H]
    \centering
    \includegraphics[width=0.82\textwidth]{chapter8_target_ranking.png}
    \caption{数值特征与对数互动量的Spearman相关系数排序}
    \label{fig:chapter8_target}
\end{figure}

粉丝数对数与对数互动量的相关性最强，Pearson相关系数为0.489，Spearman相关系数为0.469，均在FDR校正后显著。这表明发布者既有受众规模仍是笔记互动表现最重要的单一关联因素，但相关系数未达到高度相关水平，笔记热度并不能仅由粉丝基数解释。正文长度和全文长度与互动量呈弱负相关，Spearman系数分别为$-0.116$和$-0.115$；是否具有有效发布地点的系数为$-0.111$。情绪极端程度、标题正文重合度和正向情感概率的Spearman系数分别为$-0.070$、$-0.068$和$-0.063$，方向与第3章的回归结果基本一致，但效应强度较小。标签数量、表情符号、星期、感叹号和问号等变量的绝对相关系数均小于0.05，统计显著性主要受大样本量影响，不宜解释为具有较强的实际作用。

为进一步比较高热度内容，本文将总互动量位于样本前10\%的笔记定义为“爆款”，并对连续特征实施点二列相关、Welch独立样本$t$检验和Mann--Whitney U检验，同时报告Cohen's $d$。粉丝数对数的组间差异最明显：普通组均值为7.613，爆款组均值为9.589，$d=0.696$，属于中等偏大的效应。相比之下，爆款组正文长度和全文长度更短，对应效应量分别为$-0.213$和$-0.211$；标签数量、情绪极端程度与表情数量的效应量绝对值约为0.12至0.14，其余变量的效应量更小。问号数量和星期编码在FDR校正后均未表现出稳定的爆款组差异。

星期编码本身不适合作为严格连续变量解释，因此本文另行进行分组描述。如图\ref{fig:chapter8_weekday}所示，周六和周日的爆款比例分别为10.91\%和10.86\%，周一为10.42\%，周五最低，为8.54\%。该结果说明星期差异主要体现为类别间波动，而非随编码增加或减少的线性趋势。

\begin{figure}[H]
    \centering
    \includegraphics[width=0.9\textwidth]{chapter8_weekday.png}
    \caption{不同星期的互动量分布与爆款比例}
    \label{fig:chapter8_weekday}
\end{figure}

总体而言，粉丝规模是与互动量及爆款状态关联最稳定的特征；文本长度表现出弱负相关；情感、标签及格式变量虽然在大样本下多达到统计显著，但实际效应普遍较弱。因此，后续预测模型不应仅根据$p$值筛选变量，而应综合考虑效应量、非线性结构与变量的理论意义。

\section{非线性关系探索}
线性相关系数可能掩盖阈值、平台期或转折关系。为此，本文对粉丝数对数、全文长度、标签数量、正向情感概率、情绪极端程度和标题正文重合度进行LOWESS局部加权平滑，并将全部有效样本划分为20个分位数组计算组内平均互动。为减少极端值对曲线边界的影响，作图时剔除各自变量上下0.5\%的观测；散点和LOWESS最多随机抽取20,000条记录，而分位数组均值仍使用全部有效数据。

图\ref{fig:chapter8_lowess_main}展示了粉丝规模与文本长度的主要非线性形态。粉丝数对数与互动量总体呈明显正向关系，但曲线并非全程保持相同斜率，说明粉丝增长带来的互动增量存在阶段差异。全文长度曲线则表现出较明显的弯曲：短文本区间互动水平相对较高，随着长度增加先下降，随后在部分较长文本区间趋于平缓或出现局部回升。这与第3.1节在OLS中加入长度平方项后得到的非线性结论相互印证。

\begin{figure}[H]
    \centering
    \begin{minipage}{0.49\textwidth}
        \centering
        \includegraphics[width=\textwidth]{chapter8_lowess_fans.png}
    \end{minipage}
    \begin{minipage}{0.49\textwidth}
        \centering
        \includegraphics[width=\textwidth]{chapter8_lowess_content.png}
    \end{minipage}
    \caption{粉丝规模与全文长度的LOWESS非线性关系}
    \label{fig:chapter8_lowess_main}
\end{figure}

情绪极端程度的LOWESS结果见图\ref{fig:chapter8_lowess_emotion}。总体趋势显示，越偏离中性区间，平均对数互动量越低，但局部区间存在波动。这说明情感效应不宜简单表述为严格线性下降，更适合在后续模型中使用二次项、限制性立方样条或树模型检验。标签数量、正向情感概率和标题正文重合度也呈现不同程度的局部弯曲，但其总体变化幅度明显小于粉丝规模和文本长度。

\begin{figure}[H]
    \centering
    \includegraphics[width=0.78\textwidth]{chapter8_lowess_emotion.png}
    \caption{情绪极端程度与互动量的LOWESS关系}
    \label{fig:chapter8_lowess_emotion}
\end{figure}

需要指出，LOWESS属于探索性可视化方法，曲线形态不能直接证明因果关系。账号粉丝基础、内容类别、采集时间以及平台推荐机制都可能同时影响文本特征和互动结果。因此，本节的非线性发现主要用于为后续模型设定提供依据，而不将曲线转折点直接解释为普遍适用的内容运营最优值。

\chapterimage{chapter_head_2.pdf}
\chapter{特征选择与重要性评估}

\section{基于统计检验的初步筛选}
在构建复杂的机器学习模型之前，通过统计检验对大量原始特征进行初步筛选 , 可以剔除与目标变量无关或相关性极弱的特征， 降低维度、减少噪声，并提高后续建模的效率与可解释性 。本部分将根据任务类型（回归或分类）采用不同的统计检验方法。
\begin{itemize}
	\item \textbf{针对回归任务（预测互动量）：} 数值型特征：对每个数值特征x，计算其与目标变量y（通常取对数化后的点赞数） 的皮尔逊相关系数r及显著性p值。若p < 0.05，则认为该特征与目标存在显著线性相关，予以保留；否则剔除。由于皮尔逊相关对线性关系敏感，对于可能呈非线性关系的特征（如标签数量） ，可同时计算斯皮尔曼秩相关系数作为补充。
分类型特征：对于类别型特征，使用单因素方差分析（ANOVA） 检验不同类别下目标变量均值是否存在显著差异 。若F检验的p < 0.05 ，则保留该特征；否则剔除。对于类别数较多的特征，可先合并低频类别后再进行检验。
	\item \textbf{针对分类任务（预测是否爆款）：} 数值型特征：将样本按爆款标签分为两组 ，首先进行正态性检验（Shapiro- Wilk或Kolmogorov-Smirnov） 。若两组数据均近似正态分布，则使用独立样本t检验比较均值差异；若任一组的分布非正态 ，则使用Mann-Whitney U检验比较中位数差异。保留p < 0.05的特征。
分类型特征：构建特征与爆款标签的列联表 ，进行卡方独立性检验 。若p < 0.05，则拒绝原假设（特征与爆款独立），保留该特征。对于有序分类特征（如星期几），可同时进行趋势卡方检验。
\end{itemize}

\section{基于机器学习模型的特征重要性评估}
统计检验主要捕捉特征与目标之间的线性或单调关系，而机器学习模型能够自动学习复杂的非线性模式及特征间的交互作用。通过多种模型的内部机制，我们可以获得更全面、更稳健的特征重要性排序，为最终构建可解释的预测模型和提出优化策略提供依据。

\begin{itemize}
	\item \textbf{树模型的重要性：} 随机森林（Random Forest）：训练回归或分类随机森林模型，获取基于基尼不纯度减少（或均方误差减少）的特征重要性评分。该评分衡量了每个特征在树节点分裂时减少不纯度的累计贡献。通过交叉验证选择树的数量和最大深度等超参数；梯度提升树：训练XGBoost或LightGBM模型，获取基于分裂增益（gain）的特征重要性，即每个特征在每次分裂时带来的平均损失减少量。该指标比基于分裂次数的权重更能反映特征的预测能力。
	\item \textbf{正则化线性模型：} Lasso回归（L1正则化）：对特征进行标准化后，训练Lasso回归模型。Lasso通过L1惩罚项将部分特征的系数压缩为零，从而实现嵌入式特征选择。保留系数非零的特征，其绝对值大小反映了特征对预测的贡献程度。由于Lasso对线性关系敏感，可与树模型的结果互补。
	\item \textbf{模型不可知解释方法：} SHAP值：基于博弈论中的Shapley值，为每个样本的每个特征分配一个贡献值，表示该特征对预测结果的边际贡献。计算所有样本的SHAP值绝对值平均值，得到全局特征重要性排序。SHAP值不仅能反映重要性，还能通过依赖图展示特征与预测值的关方向及非线性模式。对于XGBoost或随机森林模型，可使用TreeExplainer高效计算SHAP值。
\end{itemize}


将上述三种方法得到的重要性评分进行归一化处理，并采用加权平均或排名平均的方法获得综合重要性得分。设定阈值，筛选出核心特征集。同时，观察不同方法的一致性：若某个特征在多种方法中均排名靠前，则其重要性更为可信。
